\documentclass[12pt, a4paper]{article}

% PACKAGES
\usepackage[utf8]{inputenc}
\usepackage{amsmath}
\usepackage{amssymb}
\usepackage[margin=1.2in]{geometry}
\usepackage{authblk}
\usepackage{bm}
\linespread{1.2}

\title{A Reflection on the Electron: The Standing Wave Mechanisms and Spin Topology}
\author{Haifeng Qiu}

\begin{document}

\begin{center}
    {\uppercase{\Large \textbf{A Reflection on the Electron}}} \\
    \vspace{0.5em}
    {\large \textit{The Standing Wave Mechanisms and Spin Topology}} \\
    \vspace{2em}
    \textbf{Haifeng Qiu} \\
    \today
    \vspace{2em}
\end{center}

\begin{abstract}
In the previous article, ``A Reflection on the Magnetic Field,'' we elucidated the physical ontology of electromagnetic fields and waves as microscopic ``Moiré patterns'' on the computational grid. This paper extends this logical chain further into the microscopic realm, striking directly at the most fundamental unit of matter---the electron. Starting from the underlying propagation conditions of electromagnetic waves, we re-examine the true physical entity morphology of the electron within the atom, attempting to reveal the profound grid dynamical mechanisms behind its spin and magnetic moment. We propose that the electron is not a ``point particle'' in the traditional sense, but rather a macroscopic electromagnetic standing wave formed by ``entropy momentum density fluctuations'' (photons) on the underlying computational grid, which undergo path bending and self-trapping under their self-generated ``momentum pressure.'' By introducing the ``beat delay'' mechanism of cross-layer information and the environmental dependence of the effective speed of light, this model provides a purely mechanistic and grid-dynamical explanation for the ontology of charge, the equivalence of mass and charge, the radiation mechanism of quantum transitions, and the abstract $4\pi$ topological symmetry of spin.
Furthermore, through rigorous mathematical modeling using higher-order localized equations and infinite-order non-local smoothing operators, we resolve the point-charge singularity and mathematically unify the exponential density source with a Gaussian-like effective potential field. Then, by mapping the static grid scale to dynamic fluid transport coefficients, we demonstrate that the Schrödinger equation can be a macroscopic linear emergence of the underlying fluid's continuity equation.
\end{abstract}

\vspace{2em}

\section{Theoretical Reconstruction of Electromagnetic Wave Propagation Conditions}

Electromagnetic waves are essentially the propagation of accelerated changes in entropy momentum density (corresponding to the classical magnetic vector potential) through space. As a photon propagates, it constantly generates electric field momentum in the direction perpendicular to its propagation. To accurately describe the interactions between different charges, we must stratify and refine the electric field effects at the $b_0-b_1$ layer, and deepen and correct the initial assumptions regarding ``entropy momentum'' in the theory.

\subsection{The Microscopic Essence of Entropy Momentum: ``Beat Delay'' and the Ontology of Charge}

\begin{itemize}
    \item \textbf{Cross-Layer Beat Delay:} Entropy momentum itself is a microscopic bit-timing pattern. We consider it to be a cross-computational-layer \textbf{beat delay mechanism}, manifested as a deterministic timing relationship between the $b_0$ and $b_1$ layers. That is, a pattern first appears on $b_0$ and is then copied to $b_1$, or it first appears on $b_1$ and is then copied to $b_0$.
    \item \textbf{Generation of Electric Field Momentum:} This intrinsic timing asymmetry cannot spontaneously reach equilibrium. Its propagation in different directions through space constitutes the \textbf{physical ontology of electric field momentum}.
    \item \textbf{Definition of Positive/Negative Charge:} The two opposite beat timings described above correspond respectively to the fields generated by two different charges. We define these as \textbf{positive charge momentum} and \textbf{negative charge momentum} (abbreviated as positive momentum and negative momentum). A charge is essentially a generator that continuously creates this ``beat difference'' between these layers.
\end{itemize}

\subsection{The Differentiation of Photons and the Emergence of Absolute Light Speed}

\begin{itemize}
    \item \textbf{Positive Photons and Negative Photons:} Due to the different momentum generation mechanisms, the microscopic morphologies of the entropy momentum density fluctuations generated by positive and negative charges are vastly different. We distinguish them as \textbf{positive photons} and \textbf{negative photons}. Because of the spatiotemporal asymmetry of their massive microscopic information structures, two beams of opposite-type photons will not cancel each other out when they meet.
    \item \textbf{Macroscopic Equivalence:} Although their microscopic structures differ, if we only consider spatial relativistic effects, the essence of a photon is a wave; fluctuations of thrust and pull are equivalent in their effects, and the two can transform into each other during interactions with atomic nuclei. Meanwhile, the propagation and interaction of entropy momentum within the computational layer strictly follow the \textbf{principle of linear superposition}.
    \item \textbf{Environmental Dependence of Light Speed (Momentum Pressure Effect):} The speed of light is an emergent phenomenon of the underlying grid. The accelerated variation of electric field momentum density propagates in the perpendicular direction. \textbf{Its effective propagation speed is constrained by the ``momentum pressure'' of the space it occupies}.
    \begin{itemize}
        \item \textbf{Homogeneous Deceleration:} Pre-existing negative momentum in space exerts a ``pressure'' on a negative photon propagating perpendicular to the momentum direction. This rigid pressure constitutes a propagation resistance, thereby exerting a \textbf{decelerating effect} on the negative photon. Furthermore, the effects of the same type of momentum moving toward each other will not cancel out, but will instead be strengthened.
        \item \textbf{Heterogeneous Acceleration:} Conversely, positive momentum in the environment reduces the propagation resistance for negative photons, \textbf{accelerating} their motion, though its upper limit remains constrained by the speed of light in a vacuum.
    \end{itemize}
\end{itemize}

\section{The Ontology of the Electron: A Self-Trapped Photon Standing Wave}

Within the framework of \textit{``Computational Realism,''} \textbf{the physical entity of the electron cloud (probability cloud) itself is the entirety of the electron.}

\begin{itemize}
    \item \textbf{The Electron Entity as an Electromagnetic Standing Wave:} The electron is a macroscopic standing wave system composed of electromagnetic waves (photon groups) moving in circular motion (negative charges are composed of negative photons, positive charges of positive photons). The continuous existence of photons creates electric field entropy momentum, and the electric field entropy momentum in turn distorts the propagation paths of the photons. The two act as cause and effect for each other, maintaining an eternal cycle of photon entropy momentum.
    \item \textbf{Self-Trapping Mechanism Induced by ``Refractive Index'':} Within the spatial region of the electron cloud, the circulating photons continuously ``pump out'' a fixed entropy momentum. Although the electric field vector intensity at the geometric center of the electron cloud is zero (consistent with QED models), its \textbf{electric field pressure density (a scalar) is not zero at the center; rather, it rises sharply the closer it gets to the center.}
    \item \textbf{Dynamics of Path Bending:} According to the aforementioned ``momentum pressure'' deceleration effect, the closer to the center, the lower the effective propagation speed of the photon. This extreme spatial gradient of the speed of light forms a powerful \textbf{``refractive index''} on the grid, forcibly twisting the linear motion paths of the photons and compelling them into circular motion. However, the photon density does not increase infinitely; as its density increases, photons will gradually deflect from tangential motion toward the center. Due to geometric relationships, the central momentum pressure is always greater than that on the outside. Eventually, photons will penetrate the center and propagate outward again; this process limits the photon density.
\end{itemize}

\section{The Scale Morphology and Dynamical Transitions of the Electron}

The electron is a flexible dynamical structure deeply coupled with its environment. Its volume and morphology depend on the standing wave conditions achieved by the interaction between photons and the electric field.

\begin{itemize}
    \item \textbf{The Boundary Line of Volume (Speed of Light Limit Constraint):} On the outer edge of the electron cloud, as the diameter increases, the electric field pressure density drops rapidly, and the speed of the outer circulating photons increases accordingly. When the peripheral speed of light approaches the absolute vacuum speed of light $c$, the speed encounters the physical bottleneck of the grid and cannot increase further. This physical limit destroys the proportional requirements for maintaining the standing wave, making it impossible for the volume of the electron entity to expand infinitely. \textbf{This limit boundary macroscopically determines the exact magnitude of the elementary charge $e$.}
    \item \textbf{Equivalence of Electron Mass and Charge:} The mass of an electron is always perfectly equivalent to its charge. Ontologically, both represent the same fact: exactly how many circulating photons are ``bound'' within this standing wave structure (these photons pump out momentum at a constant rate).
    \item \textbf{Environmental Dependence and Photon Emission (The Essence of Quantum Transitions):} The morphology of the electron cloud always represents the effective standing wave form that can be maintained under the current environment. 
    \begin{itemize}
        \item \textbf{Mechanism of Energy Level Transitions:} When a free electron falls into the potential well of an atomic nucleus, due to the accelerating effect of the nucleus's positive momentum environment on the electron's negative photon flow, the \textbf{effective speed of light} of the photons inside the electron \textbf{increases}. This directly leads to the collapse of the original standing wave proportional conditions; the outer structures of the electron cloud are stripped away, and the excess energy is radiated outward in the form of \textbf{free photons}.
        \item Correspondingly, in the electron orbitals outside the nucleus, due to geometric relationships, the closer a region is to the nucleus, the greater the spatial rate of change of the nucleus's positive electric field momentum. This causes the curvature of the electron cloud to become greater, thereby shaping the morphologies of electron clouds at different energy levels.
    \end{itemize}
\end{itemize}

\section{The Topological Code of Angular Momentum and Magnetic Moment}

The abstract properties of spin and magnetic moment in the Standard Model possess intuitive explanations rooted in pure geometry and fluid topology under this dynamical framework.

\begin{itemize}
    \item \textbf{The Dynamic Origin of Magnetic Moment:} Because the photons inside the electron cloud undergo high-speed periodic motion, the entropy momentum density they pump outward is not absolutely constant in its spatiotemporal distribution. This high-frequency periodic fluctuation in density directly constitutes the source of the macroscopic magnetic moment.
    \item \textbf{Spin Angular Momentum and Orbital Angular Momentum:} Essentially, these are an effective physical decomposition of the same macroscopic entity. The reality is that the topological phase cycles of the photons are acted upon by external electric field entropy momentum, causing a reorganization of their shape and scale.
    \item \textbf{The Topological Essence of Spin-1/2 (Dual Orthogonal Precession):} Within the entropy momentum phase of the electron standing wave, there actually exist \textbf{two mutually orthogonal, independent ``precessional'' cycles} (geometrically analogous to satellite orbits with a precession axis)\footnote{This description is not entirely accurate; a more precise physical picture is provided in Section 12.}. To maintain the stability of the standing wave structure, the cycle periods of these two dimensions must remain absolutely consistent.
    \item \textbf{Constraints of the Dual Cycle and Lorentz Effects:} When an electron is acted upon by external electric field entropy momentum, differences in effective propagation speeds arise across different regions. To strictly satisfy the consistency requirement for the circulation speeds in these two orthogonal directions, the overall topological morphology of the electron must reorganize. This requirement for internal periodic synchronization is the root cause of the electron's length contraction effect. Furthermore, when the electron moves, the actual grid motion path of the photons becomes larger, allowing the standing wave mode to accommodate more mass.
    \item \textbf{Geometric Decoding of $4\pi$ Symmetry:} The fundamental reason why a spin-1/2 angular momentum must rotate $720^\circ$ ($4\pi$) in space to return to its original state is that the phase standing wave morphology of its entropy momentum density is based on a dual orthogonal cycle on a spherical surface. A rotation of $360^\circ$ ($2\pi$) in only a single plane cannot simultaneously reset the cycle of the other orthogonal phase.
    \item \textbf{The Dynamical Essence of the Pauli Exclusion Principle: Topological Meshing of Complementary Phases} \\
    That an electron is a standing wave of photons means it must preserve the fluctuating state of momentum density. Two electrons with identical spins (co-directional $4\pi$ chiral precession) cannot coexist by meshing their momentum densities. However, when two electrons possess ``opposite spins,'' the two standing wave systems serve as \textbf{perfect mirrors} of each other across their dual orthogonal topologies. When they occupy the same orbital space, the dense and sparse regions of their photon momentum flows complement each other, and their directions of motion are opposite. This ``perfect phase complementarity'' causes the two high-frequency photon streams to act like perfectly engaged four-dimensional mechanical gears, interpenetrating in three-dimensional space without mutual interference. They share the same macroscopic volume, but their microscopic photon momentum flows are perfectly staggered in four-dimensional spatiotemporal coordinates, thereby avoiding the superimposition of ``momentum pressure.''
    
    This topological meshing directly triggers two decisive macroscopic physical consequences:
    \begin{enumerate}
        \item \textbf{Magnetic Moment Returns to Zero (Texture Dissolution):} Because the moving dense and sparse regions are perfectly leveled out by mirroring, the ``microscopic Moiré patterns'' bearing periodic fluctuations, which originally radiated outward individually, are internally canceled out at extremely close range. The system no longer leaks dynamic momentum texture residues to the outside, which means \textbf{the macroscopic magnetic moment of the composite system is flattened to 0}.
        \item \textbf{Doubling of the Electrostatic Field:} Although the dynamic textures are canceled out, this does not erase the foundational computational fact that the two electrons serve as underlying ``beat generators.'' They continue to independently pump out negative momentum flows, thereby causing the outwardly apparent electrostatic pressure to exactly double ($-2e$).
    \end{enumerate}
\end{itemize}


\section{The Dynamics of Scalar Pressure}

In classical electromagnetism, the electric potential (a scalar) is merely the spatial integral of the field strength (a vector); the two are distinct mathematical representations of the same mathematical entity. We attempt to rethink its physical essence through the lens of the entropy momentum vector sum and entropy momentum pressure.

\subsection{The Ontology of Entropy Momentum and Electric Potential}

When a charge (the beat generator) continuously pumps out entropy momentum within the vacuum grid, the physical state of any given point in space is determined by two distinctly different parameters:

\begin{enumerate}
    \item \textbf{Field Strength (Vector Sum $\Sigma\vec{p}$):} Represents the net direction and net scouring force of the entropy momentum flow at that point. When the system is situated at a special geometric midpoint of high multi-body symmetry, the effects of the vector sum may perfectly cancel out to zero.
    \item \textbf{Momentum Pressure (Scalar Sum $\Sigma|\vec{p}|$):} Represents the ``crowdedness'' or ``vacuum strain energy density'' of the underlying grid saturated with information momentum at that point.
\end{enumerate}

We propose the following hypothesis: \textbf{The ``electric potential ($V$)'' in macroscopic classical electromagnetism finds its microscopic physical equivalent in the ``entropy momentum scalar pressure'' upon the computational grid.}

This scalar pressure possesses properties of additivity and positive/negative cancellation. The convergence of momentum from identical charges leads to a sharp surge in local grid pressure (creating a positive or negative potential barrier). Conversely, when positive and negative momenta converge, their inter-layer information phases complement each other, mutually neutralizing and causing a drop in the total system pressure in that region (the release of vacuum tension). The electric potential is a physical entity, representing the degree of congestion on the computational grid. Field strength, on the other hand, is the vector sum of the electric field's entropy momentum, representing the net exerted effect.

That the electric potential equals the path integral of the field strength is profoundly connected to spatiotemporal geometry. Because charges are point-like and the electric field diverges spherically, momentum pressure and the momentum vector sum form an algebraic identity in space. This is both the cause and the consequence of energy conservation at the charge level: the point-like divergence of a charge is the reason for the conservation of its electric potential energy and electric field force energy, while this point-like divergence is simultaneously the inevitable result of the charge's constituents (photons) being subjected to deeper laws of energy conservation. If there existed a planar, non-divergent electric field source (such as a hypothetical planar charge), the momentum pressure and momentum vector could not be converted via path integration, and the law of energy conservation would consequently fail under such a premise.

\subsection{The Essence of Electromagnetic Force: Mass's Response to the Environment and the ``Relaxation Pathway''}

Let us consider the oldest question in physics: \textbf{Why is a negative charge ``repelled by a force'' within a homologous negative momentum field?}

Within our framework, force is not an action-at-a-distance push or pull, but rather \textbf{a ``geometric adaptation and translation'' of mass (the standing wave) in response to a non-uniform environment to maintain its internal structure.}

\begin{enumerate}
    \item \textbf{Photon Inertia and Internal Restoring Force:} The photons constituting the electron are essentially wave signals, possessing a physical inertia to maintain linear propagation. The reason they are bent into the circular standing wave of the electron is due to the forced ``trapping'' by the extremely high pressure-induced refractive index at the core. Therefore, the interior of the electron is always filled with an outward ``structural tension'' or a ``restoring force'' attempting to straighten the path.
    \item \textbf{The Refractive Index Barrier of an Asymmetric Environment:} When another negative charge approaches, it alters the distribution of scalar pressure in the vacuum. On the side facing the external negative charge, the pressure surges, and the effective speed of light slows down extremely, forming a rigid ``high-refractive-index optical wall''; whereas on the side facing away from the external negative charge, the scalar pressure is relatively lower.
    \item \textbf{Relaxation Pathway and Macroscopic Force:} At this juncture, the photon flow within the standing wave encounters this highly asymmetric pressure field during each cycle. On the side facing away from the external repelling source, the lower pressure (and faster speed of light) creates a \textbf{``possible relaxation pathway''} where the system's stress can be released. The photons' ``linear inertial restoring force'' finds a breakthrough in this direction, causing the wave's geometric trajectory to slightly expand outward at this location.
\end{enumerate}

\textbf{Conclusion:}\\
The microscopic essence of the so-called ``electrostatic repulsion'' observed macroscopically is this: Driven by the internal linear inertia of the fluid standing wave composed of photons, and much like water flowing to the lowest point, it continuously undergoes holistic phase reorganization and geometric slipping toward the ``relaxation pathway'' in space where pressure is lower and resistance is weaker. $F=ma$ is not the cause, but rather the macroscopic manifestation of a series of continuous topological displacements forced upon the standing wave in a gradient refractive index field to relieve its internal structural tension.


\section{A Simplified Model of the Free Electron}

In this theoretical model, the free electron is a macroscopic standing wave fluid entity formed by photons (fluctuations in entropy momentum density) on the underlying computational grid, constrained by environmental momentum pressure. To ensure the mathematical and mechanical self-consistency of this physical picture, this section quantitatively deduces the geometric morphology and dynamical parameters of the free electron.

\subsection{Standing Wave Conditions}

Assume the electron entity is a stable standing wave formed by photons undergoing closed circular motion in the core region. According to the Eikonal Equation in geometrical optics, the path radius of curvature $R_c$ of a light ray in a non-uniform medium satisfies the following relationship:

\begin{equation}
\frac{1}{R_c} = \frac{\hat{N} \cdot \nabla n}{n} 
\end{equation}

To maintain an ideal circular trajectory, the radius of curvature of the photon path must be strictly equal to its spatial radius ($R_c = r$), and the normal vector must point toward the center ($\hat{N} = -\hat{r}$). Substituting these yields:

\begin{equation}
\frac{1}{r} = -\frac{1}{n} \frac{dn}{dr} 
\end{equation}

Integrating the above equation, we find that the spatial refractive index $n(r)$ must be strictly inversely proportional to the radius: $n(r) = \frac{n_0}{r}$.\\
Because the effective linear propagation speed of the photon on the grid is $v(r) = \frac{c}{n(r)}$ (where $c$ is the absolute vacuum speed of light on the grid), substituting the refractive index formula yields:

\begin{equation}
v(r) = \left(\frac{c}{n_0}\right) r = \omega r 
\end{equation}

This geometric constraint indicates that to maintain a closed standing wave orbit, the microscopic photon flow must exhibit an overall rigid-body rotation characteristic with a constant angular velocity ($\omega$). Note that ``rigid body'' here is a metaphor; as described earlier, the constituent photons of the actual electron exhibit precessional orbital motion, but the angular velocities of the inner and outer layers remain consistent.

\subsection{The Relationship Between Electric Potential and Light Speed, and Spatial Distribution}

This model assumes that the effective speed of light $v$ in the grid space is modulated by the local scalar electric potential (momentum pressure) $U$.\footnote{Here, the potential refers to the electric potential associated with a specific angular velocity $\omega$ of the soliton; a detailed analysis is provided in Section 11.} In a low-order approximation, this modulation mechanism is set to follow an algebraic conservation form between the momentum flow energy and the static strain pressure:

\begin{equation}
U(v) = U_{max} - \alpha v^2 
\end{equation}

Where $U_{max}$ is the maximum potential pressure at the geometric center of the system, and $\alpha$ is the modulation constant. For ease of expression, $U(v)$ and $U_{max}$ are taken as positive values (the absolute value is taken for negative potential), and the same applies below.\\
Substituting the standing wave constraint $v(r) = \omega r$ into the aforementioned modulation law, we can deduce that the spatial potential distribution inside the free electron assumes an inverted parabolic form:

\begin{equation}
U(r) = U_{max} - \alpha \omega^2 r^2 = U_{max} - \kappa r^2 
\end{equation}

\subsection{Density Distribution of the Free Electron (Poisson's Equation)}

To determine the physical entity (charge) distribution $\rho(r)$ that generates the aforementioned inverted parabolic potential, we introduce the classical three-dimensional Poisson's equation in spherical coordinates: $\nabla^2 U = -\frac{\rho}{\epsilon_0}$.\\
Solving the Laplace operator for $U(r) = U_{max} - \kappa r^2$:

\begin{equation}
\nabla^2 U = \frac{1}{r^2} \frac{d}{dr} \left( r^2 \frac{d}{dr} (U_{max} - \kappa r^2) \right) = \frac{1}{r^2} \frac{d}{dr} (-2\kappa r^3) = -6\kappa 
\end{equation}

Substituting this into Poisson's equation yields:

\begin{equation}
\rho(r) = 6\kappa\epsilon_0 = \text{Constant} 
\end{equation}

Under the dual constraints of the inverted parabolic potential and the geometric standing wave, the free electron manifests in space as a physical entity with an absolutely uniform internal density, rather than an exponentially decaying continuous probability cloud.

\subsection{Boundary Conditions}

In this model, the spatial scale of the entity is directly limited by the computational boundary of the underlying grid (the absolute speed of light $c$).\\
According to the velocity spatial distribution $v(r) = \omega r$, as the radius increases, the effective linear speed of the internal photons increases accordingly. When extending to a specific radius $R$, the photon linear velocity reaches the absolute vacuum speed of light $c$. Because the physical grid cannot transmit information faster than light, the geometric surface $v(R)=c$ constitutes the absolute physical boundary of this fluid standing wave. The system cannot maintain self-consistent circular motion equations beyond this radius.


\section{The 1s Orbital Electron Model}

The free electron model in Section 6 merely reflects a simplified model of the fluid's low-order effects and cannot completely describe the multi-body dynamical morphology of an orbital electron in the presence of an external complex electromagnetic field (such as a nuclear potential field). To adapt this fluid standing wave model to more complex physical realities, we further elevate the physical complexity of the underlying field dynamics mechanism.

\subsection{Introduction of the Fourth-Order Biphasic Equation}

We postulate that in the high-frequency dynamic interactions of multi-body systems, due to the complex action of light speed delay (retarded potentials), the interaction between the photon momentum flow and the spatial potential at the microscopic scale no longer exhibits simple long-range Coulombic behavior, but rather presents localized shielding characteristics akin to the Yukawa potential.

To accurately describe this localized interaction process affected by light speed delay, this theoretical model introduces the Bopp-Podolsky fourth-order biphasic operator to replace the classical second-order Poisson's equation:

\begin{equation}
\nabla^2 (1 - a^2 \nabla^2) U(r) = - \frac{\rho(r)}{\epsilon_0} 
\end{equation}

Where $a$ is a characteristic parameter representing the localized shielding and microscopic smoothing scale. This equation, through the $\nabla^4$ operator, endows spatial field transmission with a microscopic phase smoothing response mechanism. Here, the charge density is the algebraic sum of positive and negative charge densities.

\subsection{Exact Analytical Solution and Its Properties}

When the source charge density adopts the exponentially decaying envelope macroscopically presented in physical reality (i.e., the Laplace distribution $\rho(r) \propto e^{-\mu r}$), substituting it into the aforementioned fourth-order biphasic field equation and utilizing the corresponding non-singular Green's function for three-dimensional spatial convolution integration, we can obtain the exact analytical closed-form solution for the potential field distribution:

\begin{equation}
U_{exact}(r) = \frac{1}{r} \left[ C_1 + C_2 e^{-\mu r} + C_3 e^{-r/a} + C_4 r e^{-\mu r} \right] 
\end{equation}

To ensure that the potential field does not diverge at the spatial geometric center of the system ($r \to 0$) and that the total energy remains finite, the system parameters must automatically satisfy the boundary constraint condition: $C_1 + C_2 + C_3 = 0$.

Let us analyze the algebraic properties of this exact solution in the microscopic core region. Performing a Taylor series expansion of $U_{exact}(r)$ near $r=0$, due to the aforementioned zero-point constraint condition and the inherent smoothing characteristics of the fourth-order operator, the $r^1$ and $r^3$ terms in the expansion are strictly eliminated mathematically:
\begin{equation}
U_{exact}(r) = U_0 - \kappa r^2 + \beta r^4 + O(r^5) 
\end{equation}
This deduction demonstrates that under the constraints of the fourth-order localized field equation, the original Laplace continuous density distribution with a sharp vertex has its effective potential field excited in the core region transformed into a series, of which the dominant second-order term $U_0 - \kappa r^2$ reproduces the fundamental inverted parabolic potential barrier structure required to maintain the photon circular standing wave.

\subsection{The Relationship and Properties of Potential-Light Speed Modulation}

To maintain the fluid topological closure of the orbital standing wave, the system must satisfy the purely kinematic rigid correlation condition $r = v/\omega$ at any radius. By directly substituting this velocity-radius mapping relationship into the aforementioned exact analytical solution of the potential $U_{exact}(r)$, we can obtain the exact closed mathematical expression for the modulation of light speed by the underlying grid potential:

\begin{equation}
U(v) = \frac{\omega}{v} \left[ C_1 + C_2 e^{-\frac{\mu v}{\omega}} + C_3 e^{-\frac{v}{a \omega}} + C_4 \left(\frac{v}{\omega}\right) e^{-\frac{\mu v}{\omega}} \right] 
\end{equation}

This analytical solution represents the full picture of the nonlinear modulation law incorporating localized delay effects. We convert this modulation relationship into a series form:
\begin{equation}
U(v) = U_0 - \left(\frac{\kappa}{\omega^2}\right) v^2 + \left(\frac{\beta}{\omega^4}\right) v^4 + O(v^5)
\end{equation}


\section{Another Orbital Electron Model: Exact Solution Based on Global Gaussian Even Function Potential}

In the derivation of Section 7, although the introduction of the fourth-order biphasic equation successfully eliminated the $r^1$ and $r^3$ terms, it remained impossible to achieve complete analytical cancellation of the fifth-order term ($r^5$). To strictly implement the physical hypotheses that the ``underlying computational grid is an isotropic medium'' and that the ``potential field is a pure even function,'' this section attempts to perform a reverse derivation of the corresponding charge density distribution starting from a globally smooth, pure even function potential field.

\subsection{Gaussian Potential Hypothesis and Kinematic Self-Consistency}

A more physically complete potential field for an orbital electron should adopt a Gaussian distribution form. This not only satisfies the requirement for smoothness at $r=0$ (no cusp) but is also mathematically a series sum of purely even functions:

\begin{equation}
U(r) = U_0 e^{-\alpha r^2}
\end{equation}

Where $U_0$ represents the maximum momentum pressure at the core. The Taylor expansion of this potential field in the core region perfectly preserves the ``inverted parabolic potential barrier'' required to maintain the photon standing wave, while simultaneously providing natural higher-order damping.

\subsection{Density Distribution Derivation Based on the Bopp-Podolsky Equation}

To determine the entity density distribution $\rho(r)$ that generates this potential field, we substitute it into the fourth-order biphasic field equation:

\begin{equation}
\nabla^2(1 - a^2\nabla^2)U(r) = -\frac{\rho(r)}{\epsilon_0}
\end{equation}

Utilizing the Laplace operator in spherical coordinates $\nabla^2 = \frac{d^2}{dr^2} + \frac{2}{r}\frac{d}{dr}$, through precise differential operations, we obtain:

\begin{enumerate}
    \item \textbf{Second-Order Term:}
    \begin{equation}
    \nabla^2 U(r) = U_0 (4\alpha^2 r^2 - 6\alpha) e^{-\alpha r^2}
    \end{equation}
    \item \textbf{Fourth-Order Term:}
    \begin{equation}
    \nabla^4 U(r) = U_0 (16\alpha^4 r^4 - 80\alpha^3 r^2 + 60\alpha^2) e^{-\alpha r^2}
    \end{equation}
\end{enumerate}

Combining the above results and substituting them into the original equation, and extracting the common factor $e^{-\alpha r^2}$, we obtain the global exact analytical solution for the entity density:

\begin{equation}
\rho(r) = \epsilon_0 U_0 \left[ (6\alpha + 60a^2\alpha^2) - (4\alpha^2 + 80a^2\alpha^3)r^2 + 16a^2\alpha^4 r^4 \right] e^{-\alpha r^2}
\end{equation}

\subsection{The Potential-Light Speed Modulation Law under Gaussian Potential and Its Series Expansion}

According to the kinematic constraint of the fluid standing wave $r = v/\omega$ (i.e., the linear velocity of a photon within the standing wave is proportional to its radius), we directly map the global Gaussian potential $U(r)$ into velocity space, obtaining the nonlinear modulation function of the local effective speed of light by the potential:

\begin{equation}
U(v) = U_0 e^{-\left( \frac{\alpha}{\omega^2} \right) v^2}
\end{equation}

To analyze its dynamical performance across different velocity domains, we expand it into a Maclaurin series:

\begin{equation}
U(v) = U_0 \left[ 1 - \frac{\alpha}{\omega^2} v^2 + \frac{1}{2} \left( \frac{\alpha}{\omega^2} \right)^2 v^4 - \frac{1}{6} \left( \frac{\alpha}{\omega^2} \right)^3 v^6 + \dots \right]
\end{equation}

Under this framework, the relationship between the potential drop and the local photon flow velocity manifests as an even-function damping series modulated by multiple higher-order terms. The ``even-power'' expansion of the potential corresponds to the ``odd-power'' response of the field strength in nonlinear optics.

Its even powers imply an isotropic synthesis rule for velocity on the spatial axes. In the sense of nonlinear optics, the underlying computational grid (vacuum) is an ``isotropic medium.'' The strict disappearance of odd-power terms in the series means that in the three-dimensional spatial coordinate system, the momentum synthesis along the $x, y, z$ axes follows a strict root-sum-square law. This ensures that when the electron is deformed by external field perturbations, the frequency synchronization of its internal standing wave will not collapse due to directional inconsistencies, thereby maintaining the topological stability of the electron entity.

Analyzing its physical properties reveals: in the low-speed core region, the second-order term dominates, and the system exhibits basic quasi-classical behavior; whereas as the local photon flow velocity $v$ increases, the weight of the nonlinear higher-order damping terms (such as $v^4, v^6$, etc.) rises rapidly. This higher-order response mechanism, derived from the fourth-order localized equation, naturally gives rise to a progressive ``Soft Asymptotic Boundary'' effect at the physical spatial boundaries, allowing the standing wave flow velocity to be smoothly and strictly confined by the absolute grid's speed of light limit $c$ through fluid damping characteristics.



\section{The Unification of Field Dynamics: Infinite-Order Smoothing Operators and Non-Local Grid Responses}

In the discussions of the previous two sections, we examined the orbital electron through two distinct pathways: Section 7 attempted to start from the exponential distribution of quantum mechanics, discovering that a fourth-order operator is sufficient to eliminate low-order singularities; Section 8 started from a perfect Gaussian even-function potential and deduced a density distribution with a shielding structure. This section will demonstrate that these are not mutually exclusive models, but rather \textbf{progressive manifestations of the same physical process under different operator orders}.

\subsection{From Bopp-Podolsky to Infinite-Order Exponential Operators}

The Bopp-Podolsky operator $(1 - a^2\nabla^2)$ introduced in Section 7 is essentially a second-order truncation of a deeper operator. According to Computational Realism, information transmission between grids is recursive. If we consider this interaction as a nearly infinite number of iterations relative to macroscopic physical phenomena, the original finite-order differential equation will evolve into a non-local operator containing \textbf{infinite-order derivatives}:

\begin{equation}
\hat{L}_{total} = \lim_{n \to \infty} \left(1 - \frac{a^2}{n}\nabla^2\right)^n = e^{-a^2\nabla^2}
\end{equation}

Here, the exponential operator $e^{a^2\nabla^2}$ is known as the \textbf{Weierstrass Transform} or the \textbf{heat kernel convolution operator}. Its essential function is to perform a global Gaussian smoothing on spatial distributions, and $e^{-a^2\nabla^2}$ is its inverse operator.

\subsection{Mathematical Unification of Exponential Distribution Sources and Gaussian Potential Fields}

Suppose the ``bare charge density'' $\rho_{bare}(r)$ of the electron entity follows an exponential distribution similar to a quantum mechanical wave function (i.e., having a core cusp $e^{-\mu r}$). When this distribution is projected through a computational grid with infinite-order smoothing properties, the resulting effective potential field $U(r)$ satisfies:

\begin{equation}
\nabla^2 \left( e^{-a^2\nabla^2} U(r) \right) = -\frac{\rho_{bare}(r)}{\epsilon_0}
\end{equation}

A simple calculation yields the relationship between the effective potential and the bare potential:

\begin{equation}
U_{eff}(r) =  e^{a^2\nabla^2} \left( \nabla^{-2} \left( -\frac{\rho_{bare}(r)}{\epsilon_0} \right) \right) =  e^{a^2\nabla^2} \left( U_{bare}(r) \right) 
\end{equation}

The relationship between the effective density and the bare density is:

\begin{equation}
\nabla^{2} \left( U_{eff}(r) \right)  =  e^{a^2\nabla^2} \left(  -\frac{\rho_{bare}(r)}{\epsilon_0} \right) = -\frac{\rho_{eff}(r)}{\epsilon_0}
\end{equation}

That is:

\begin{equation}
e^{a^2\nabla^2} \left(  \rho_{bare}(r) \right) = \rho_{eff}(r)
\end{equation}

This indicates that inside the electron, the effective potential can be viewed as the result of the bare potential diffusing in space according to the \textbf{heat kernel convolution operator}. As long as the number of grid interaction iterations is sufficiently high, an exponential-like density source $\rho \propto e^{-\mu r}$ with a core cusp can project a singularity-free Gaussian-like even-function potential field in space. Note that the relationship between the effective density and the bare density here does not imply that the bare density entity has diffused; rather, it means that the \textit{effect} of the bare density has diffused, forming an effect identical to that of the effective density.

\subsection{Physical Interpretation: The Grid as a ``Spatiotemporal Low-Pass Filter''}

\begin{itemize}
    \item \textbf{Separation of ``Bare Entity'' and ``Effective Field'':} The original momentum accumulation (bare density) of the electron on the underlying grid generates a cusp-bearing morphology similar to that in quantum mechanics. However, due to the endogenous smoothing effect of grid information transmission (nearly infinite-order recursion), the internal ``potential'' manifests as a smoothed Gaussian distribution.
    \item \textbf{The Mechanical Origin of Non-Locality:} The introduction of the infinite-order operator implies that the internal potential field no longer depends solely on the local photon density, but is rather the \textbf{result of a non-local convolution} of the charge distribution across the entire grid region. This suggests the underlying reason for the holistic coordination and non-local phase characteristics of the point charge exhibited by the electron cloud, without violating the premise of local realism.
\end{itemize}

\subsection{Conclusion: The Unification of the Electron Entity}

Through the unification in Section 9, we have established the morphology of the electron in orbit: it is a \textbf{momentum flow source characterized by an exponential distribution}, whose projected bare potential, via \textbf{infinite-order grid operators}, internally projects a \textbf{Gaussian-shaped effective field}.


\section{The Emergence of Quantum Evolution: From Fluid Topological Dynamics to the Schrödinger Equation}

This section aims to establish a precise mathematical mapping between the static topology of the underlying grid and macroscopic quantum dynamics. By converting the grid smoothing operator established in Section 9 into a time-evolution generator, we can demonstrate that a fluid standing wave constrained by the physical limits of the grid inevitably follows the Madelung continuity equation, ultimately manifesting mathematically as the linearized Schrödinger equation.

\subsection{Mapping from Static Grid Scale to Dynamic Transport Coefficients}

Section 9 proved that the electron entity exhibits a spatial smoothing effect determined by the intrinsic grid scale $a^2$: $\rho_{eff} = e^{a^2\nabla^2}\rho_{bare}$. Under the recursive interaction framework of Computational Realism, this static smoothing effect does not exist in isolation, but is the time-accumulated result of the grid exchanging information at a background angular frequency $\omega$.

At the dynamic level, this means that within a unit interaction period $T = 1/\omega$, the effective spatial smoothing area covered by the momentum flow on the grid is $a^2$. Therefore, we define the intrinsic dynamic transport coefficient (diffusivity) $D$ of this standing wave system as:

\begin{equation}
D = \frac{a^2}{T} = a^2 \omega 
\end{equation}

This equation completes the conversion from the static grid geometric parameter ($a^2$) to the dynamic fluid evolution parameter ($D$).

\subsection{The Madelung Continuity Equation Based on Standing Wave Conservation Laws}

Because the electron entity is defined as the conserved accumulation of momentum density on the underlying grid, the time evolution of its effective density $\rho$ (i.e., $\rho_{eff}$) must obey the continuity equation of classical fluid mechanics. According to Madelung's fluid dynamic formulation, if the fluid is driven rotationally and guided by a topological phase field $S$, such that its velocity field satisfies $\mathbf{v} = D \nabla S$, then the conservation law is expressed as:

\begin{equation}
\frac{\partial \rho}{\partial t} + \nabla \cdot (\rho D \nabla S) = 0 
\end{equation}

Substituting the grid docking parameter $D = a^2 \omega$ yields:

\begin{equation}
\frac{\partial \rho}{\partial t} + \nabla \cdot (\rho \cdot a^2 \omega \cdot \nabla S) = 0 
\end{equation}

At this point, the motion of the electron is no longer viewed as abstract probability migration, but rather as the genuine transport of the ``effective density'' established in Section 9 following Madelung fluid laws in space.

\subsection{Scale Anchoring and the Emergence of Physical Constants}

By combining the mass-energy relationship of the particle $E = mc^2 = \hbar\omega$ with the geometric limit condition of standing wave self-trapping $a = c/\omega$, we can precisely anchor the physical dimensions of the operator coefficients in the above equation:

\begin{equation}
a^2 \omega = \left( \frac{c}{\omega} \right)^2 \omega = \frac{c^2}{\omega} = \frac{c^2}{mc^2/\hbar} = \frac{\hbar}{m} 
\end{equation}

The derivation results show that the intrinsic smoothing scale $a$ of the grid strictly corresponds to the reduced Compton wavelength of the electron $\lambda_c = \hbar/mc$. Through this anchoring, the transport coefficient $D$ in the Madelung continuity equation is endowed with a clear physical meaning, namely $\hbar/m$.

\subsection{Linearized Dimensional Reduction and the Generation of the Schrödinger Equation}

To transform this non-linear coupled system into a computable linear system, we define a pure complex geometric construct function $\psi$ (i.e., the wave function probability amplitude), packaging the real density $\rho$ and the real phase $S$ onto the complex plane:

\begin{equation}
\psi = \sqrt{\rho} e^{iS} 
\end{equation}

\begin{equation}
\psi^* = \sqrt{\rho} e^{-iS} 
\end{equation}

Substituting the pure real current density $\mathbf{j} = \rho (\omega a^2) \nabla S$ with a form expressed using the complex numbers $\psi$ and $\psi^*$, we take the spatial gradients of the wave function and its conjugate respectively:

\begin{equation}
\nabla \psi = (\nabla \sqrt{\rho}) e^{iS} + i \sqrt{\rho} e^{iS} \nabla S 
\end{equation}

\begin{equation}
\nabla \psi^* = (\nabla \sqrt{\rho}) e^{-iS} - i \sqrt{\rho} e^{-iS} \nabla S 
\end{equation}

Substituting these back into the original continuity equation:

\begin{equation}
\frac{\partial (\psi\psi^*)}{\partial t} + \nabla \cdot \left[ -i \left(\frac{1}{2}\omega a^2\right) (\psi^* \nabla \psi - \psi \nabla \psi^*) \right] = 0 
\end{equation}

Expanding the time partial derivatives using Leibniz's rule and expanding the divergence operation (due to symmetry, the $\nabla\psi^* \cdot \nabla\psi$ terms cancel each other out):

\begin{equation}
\left( \psi^* \frac{\partial \psi}{\partial t} + \psi \frac{\partial \psi^*}{\partial t} \right) - i \left(\frac{1}{2}\omega a^2\right) \left( \psi^* \nabla^2 \psi - \psi \nabla^2 \psi^* \right) = 0 
\end{equation}

Regrouping and strictly decomposing the equation into two parts containing $\psi^*$ and $\psi$:

\begin{equation}
\psi^* \left[ \frac{\partial \psi}{\partial t} - i \left(\frac{1}{2}\omega a^2\right) \nabla^2 \psi \right] + \psi \left[ \frac{\partial \psi^*}{\partial t} + i \left(\frac{1}{2}\omega a^2\right) \nabla^2 \psi^* \right] = 0 
\end{equation}

Let the operator expression inside the square brackets be $A$:

\begin{equation}
A = \frac{\partial \psi}{\partial t} - i \left(\frac{1}{2}\omega a^2\right) \nabla^2 \psi 
\end{equation}

We obtain:

\begin{equation}
\operatorname{Re}(\psi^* A) = 0 
\end{equation}

That is, $A$ is a pure imaginary number. In the Madelung continuity equation, $A$ equals the sum of the external potential energy and the Bohm quantum potential, where the Bohm quantum potential represents photon inertia and internal restoring forces. Here, for a simplified treatment, we assume the particle is in an environment where the external potential energy and the Bohm quantum potential cancel each other out, thus $A=0$.

This perfectly splits a complex non-linear fluid equation into two decoupled, mutually conjugate complex linear partial differential equations, where the forward evolution equation is:

\begin{equation}
\frac{\partial \psi}{\partial t} = i \left(\frac{1}{2}\omega a^2\right) \nabla^2 \psi 
\end{equation}

Substituting parameters and simplifying yields:

\begin{equation}
i\hbar \frac{\partial \psi}{\partial t} = -\frac{\hbar^2}{2m} \nabla^2 \psi 
\end{equation}

\subsection{Conclusion}

The Schrödinger equation is essentially the macroscopic linear equivalent form manifested by the ``constrained grid smoothed field'' from Section 9, under the premise of following the Madelung continuity equation. Its essence is that the effective potential field (effective density) is the result of infinite-order diffusion through the underlying grid. This proves that quantum mechanics is not an axiomatic system independent of classical physics, but rather the dynamic emergence of underlying standing wave fluids at a specific scale.



\section{Covariant Spacetime Dynamics and Non-Local Field Equations}

\subsection{Four-Dimensional Covariant Smoothing Operator and Topological Closure Conditions (The Klein-Gordon Equation)}

In the actual underlying physical grid, information propagation is not only smoothed in three-dimensional space but is also covariant in four-dimensional spacetime. Therefore, we upgrade the static spatial heat kernel convolution operator from Section 9 to an \textbf{infinite-order, Lorentz-covariant smoothing operator in four-dimensional spacetime}:

\begin{equation}
\hat{K}_{4D} = e^{a^2\Box}
\end{equation}

Where $\Box = \partial_\mu \partial^\mu = \frac{1}{c^2}\frac{\partial^2}{\partial t^2} - \nabla^2$ is the D'Alembert operator.

According to the core physical hypothesis of this paper, the electron entity is a \textbf{``self-trapped standing wave.''} 
From the perspective of four-dimensional spacetime, a stably existing ``standing wave entity'' (whose rest mass does not dissipate over time) implies that under the smoothing iterations of the four-dimensional grid, it must reach an \textbf{intrinsic dynamic topological equilibrium state (scale invariance)}. 
In other words, the macroscopically stable effective fluid field $\psi$ must be an eigenstate of this grid characteristic scale operator:

\begin{equation}
a^2 \Box \psi = -\psi 
\end{equation}

Rearranging the above equation yields:

\begin{equation}
\left(\Box + \frac{1}{a^2}\right)\psi = 0
\end{equation}

We once again introduce the crucial ``physical constant anchoring'' conclusion from Section 10.3, namely that the intrinsic smoothing scale $a$ of the grid strictly corresponds to the reduced Compton wavelength of the electron: \textbf{$a = \frac{\hbar}{mc}$}. Substituting this into the above equation yields:

\begin{equation}
\left(\Box + \frac{m^2c^2}{\hbar^2}\right)\psi = 0 \quad \text{or written as} \quad \left(\partial_\mu \partial^\mu + \left(\frac{mc}{\hbar}\right)^2\right)\psi = 0
\end{equation}

This is the fundamental equation of relativistic quantum mechanics—the Klein-Gordon Equation. 
It proves that, at the scale of the four-dimensional grid, the rest mass $m$ of matter is essentially the reciprocal of the geometric radius of curvature when the four-dimensional spacetime standing wave achieves dynamic closure.

Meanwhile, according to the premise that two mutually orthogonal, independent precessional cycles actually exist inside the electron, the underlying fluid phase field cannot be described by a simple scalar $\psi$; it must encapsulate an intrinsic orthogonal directionality across four dimensions. 

To mathematically describe this ``orthogonal gear meshing on the four-dimensional grid,'' we perform a geometric manifold decomposition (Clifford Algebra) on the scalar D'Alembert operator $\Box$, and by factoring the equation, we can derive the Dirac Equation.

\subsection{Grid Dynamics and Non-Local Evolution in Four-Dimensional Spacetime}

In the preceding discussions, we established the D'Alembert operator as the spacetime smoothing operator for the four-dimensional electromagnetic potential. Here, we postulate further that the spacetime dynamics of the grid itself are driven by the D'Alembert operator.

\subsubsection{Local Action of the Discrete Grid and Macroscopic Non-Local Emergence}

We propose that the actual underlying vacuum is not a continuous manifold, but rather consists of infinitesimally small, discrete, and finite grids. In this discrete spacetime, the single interaction transfer of information is strictly confined between adjacent grids. Its most fundamental single local evolution operator can be expressed as $(1 + \epsilon^2\Box)$, where $\epsilon$ is the intrinsic characteristic scale of the grid action, and $\Box = \frac{1}{c^2}\partial_t^2 - \nabla^2$ is the D'Alembert operator.

Because macroscopic physical entities are not formed by a single iteration, but are achieved through countless underlying computational operations, we obtain:

\begin{equation}
\hat{K}_{macro} = \lim_{n\to\infty}  (1 + \epsilon^2 \Box)^n = e^{ n \epsilon^2 \Box} = e^{a^2 \Box} 
\end{equation}

We calculate a complete evolution, i.e., the number of grid cycles required for the soliton to achieve topological cycle closure, as $n = \left(\frac {2 \pi / \omega} { 2 \pi \epsilon / c}\right)^2$. Substituting this yields the characteristic scale coefficient $a = \frac{c}{\omega}$.

This operator indicates that the spacetime distribution of the electromagnetic potential is subjected to unceasing ``smoothing'' by this operator, which constitutes the fundamental dynamical cause of the \textbf{dispersion} of electric field momentum in spacetime. The non-locality of the exponential operator guarantees the holistic ``particle nature'' of the soliton, and the strong interaction between the soliton and external modes already occurs at the Compton wavelength. It also serves as a cross-scale mapping factor, rendering the shape of the bare entity negligible in the coupling between the bare entity and the effective potential field.

Here, we argue that the D'Alembert operator remains a grid emergence. The underlying discrete grid rules might be a discrete local convolution operator that is equivalent to the D'Alembert operator and preserves momentum and energy conservation.

\subsubsection{Grid Saturation Limit and Extremum Principle}

Because the grid is discrete and its information capacity is finite, there inevitably exists a \textbf{``grid saturation limit.''} This implies an absolute upper bound to the accumulation of electric field momentum (pressure) on the grid. As the density of the electric field momentum on the grid rises, its capacity to transmit electric field momentum decreases, manifesting as a retardation effect.

\subsection{Energy Density Modulation of the Speed of Light}

In the preceding text, we proposed a scalar modulation function for the speed of light. However, it must be clarified that the modulation relationship between the speed of light and the electric potential holds true under the premise of a fixed angular velocity within the soliton. That is, the speed of light is essentially not directly modulated by the absolute value of the electric potential, but rather by the \textbf{energy density} of the underlying grid. Inside the soliton, since the angular velocity $\omega$ is a constant fixed by the topological closure condition, the squared distribution of the effective potential $A^2$ is equivalent to the distribution of energy density. The retardation of the speed of light is a grid information processing ``delay'' effect triggered by high energy density.

Within the framework of ``Computational Realism,'' the electromagnetic potential represents the distribution morphology of information. Only when the potential field circulates at a specific angular velocity $\omega$ (i.e., the self-trapped standing wave state) does it exert a sustained momentum pressure on the grid. To reflect this dynamic feedback, we introduce a scalar mapping of the four-dimensional energy-momentum tensor to represent the energy density:

\begin{equation}
u_{inv} = T_{\mu\nu} U^\mu U^\nu \approx \epsilon_0 \omega^2 A_\nu A^\nu 
\end{equation}

Where $U^\mu$ is the four-velocity of the grid's rest frame of reference. 
The reason we cannot use $u_{inv} = -\frac{1}{4} F_{\mu\nu}F^{\mu\nu}$ to represent energy density here is that, according to our axioms, the grid momentum pressure density is a physical reality; the electric field strength is zero at the geometric center of the grid, yet the momentum pressure density is at its maximum there. 
This essentially embodies our Lorentz aether worldview: an absolute stationary grid exists, and the fundamental source of gauge variance is the dynamical tolerance of the grid rules towards soliton solutions. Furthermore, measurement, also being a grid solution, is inevitably gauge covariant. However, when studying the interior of a soliton, this inevitably breaks down. Expressed in condensed matter terminology, a ``symmetry breaking'' of this gauge variance of measurement occurs.

\subsection{Nonlinear Field Equation Based on Energy Density Saturation}

Based on the aforementioned grid dynamics, we can reconstruct the field equations. In an absolute vacuum (with zero momentum pressure), the equation degenerates into the classical Maxwell free-field equation $\Box A^\mu = 0$.

When high energy density accumulates, the local grid becomes ``viscous'' due to information overload, causing a drop in the effective speed of light. This nonlinear feedback loop induces self-focusing of the momentum flow. 

Analyzing the saturation characteristics of the discrete grid, we define this retardation effect as a logarithmic function $g(u)$ based on energy density.

\begin{equation}
g(u_{inv})= \kappa \ln\left(1 + \frac{u_{inv}}{u_{th}}\right)
\end{equation}

Combining this with the infinite-order evolution operator, we obtain the underlying \textbf{four-dimensional non-local nonlinear field equation}:

\begin{equation}
e^{a^2 \Box} A^\mu = \left[ 1 + \kappa \ln\left(1 + \frac{\epsilon_0 \omega^2 A_\nu A^\nu}{u_{th}}\right) \right] A^\mu 
\end{equation}

Where:

\begin{itemize}
    \item \textbf{$\kappa$}: A dimensionless constant reflecting the coupling strength of the grid.
    \item \textbf{$u_{th}$}: The \textbf{nonlinear characteristic threshold} of the vacuum energy density (with dimensions of $\text{J/m}^3$). It corresponds to the physical critical point where the underlying computational grid transitions from linear transmission to nonlinear self-trapping, and its order of magnitude is correlated with the Schwinger limit energy density.
    \item \textbf{$\omega$}: The intrinsic circulation frequency of the entity. For a single photon or an electron, this parameter is adaptively locked by its mass or wavelength.
    \item \textbf{$A^\mu$}: The four-potential, representing the grid momentum density and vector momentum flux.
\end{itemize}

\subsection{Electrostatic Field and Photon Solutions of the Field Equation}

To verify the self-consistency of this underlying field equation, we substitute the static electron core and the dynamic photon as test solutions into this nonlinear equation, respectively.

\subsubsection{Electrostatic Field Solution (Self-Shielding Mechanism of Mass Entities)}

For a static electron core, let the angular velocity $\omega$ be its intrinsic rotational frequency. At this point, the D'Alembert operator simplifies to $-\nabla^2$. Substituting into the first-order approximation of the field equation yields:

\begin{equation}
a^2 \nabla^2 \phi = -g(u) \phi 
\end{equation}

In the near-field core region, as the energy density $u$ rises sharply, the logarithmic term $g(u)$ activates rapidly. This nonlinear pressure, arising from ``grid saturation,'' acts as a virtual ``negative source,'' forcibly suppressing the infinite divergence of the electric potential. This proves that the finitude of electron mass originates from a \textbf{logarithmic self-shielding effect} of the underlying grid towards energy density, eliminating the point-charge singularity found in classical electrodynamics.

\subsubsection{Photon Solution (Spacetime Resonant Soliton)}

A photon is essentially a field envelope that diffuses laterally while its momentum density propagates forward. We assume the photon is an envelope solution translating along the $x$-axis: $A^\mu = \epsilon^\mu \Psi(x-ct, y, z)$, which is substituted into the operator. Due to the translation at the speed of light, the longitudinal second-order derivatives of time and space cancel out exactly ($\frac{1}{c^2}\partial_t^2 - \partial_x^2 = 0$), leaving the operator with only transverse (in the $y,z$ cross-section) diffusion: $\Box = -\nabla_\perp^2$. \\
Substituting and expanding the field equation yields:

\begin{equation}
a^2 \nabla_\perp^2 \Psi + g(\Psi^2) \Psi = 0 
\end{equation}

\textbf{Verification Result:} This result is exactly the \textbf{``Spatial Soliton Equation''} in nonlinear optics. The first term on the left side of the equation represents the natural tendency for transverse diffusive divergence of the photon caused by the grid, while the second term represents the light-speed retardation and nonlinear self-focusing tendency triggered by local high momentum pressure. The strict cancellation of the two allows the nonlinear feedback mechanism of the non-local grid equation to naturally manifest as a dynamic spacetime standing wave entity.

\subsection{Scale Analysis of the Photon Solution: Logarithmic Scaling Violation and Wavelength Locking}

In this model, the geometric transverse radius $R$ of the photon is dually constrained by the wavelength $\lambda$ and the energy density.

\vspace{0.5em}
\noindent \textbf{1. Correlation between Energy and Frequency:} \\
The quantum energy of a single photon is known to be $E = \hbar\omega$. According to the classical volume integral of energy, $W = \int u \, dV \approx (\epsilon_0 \omega^2 \Psi^2) \cdot (R^2 \lambda)$. \\
To satisfy $W = \hbar\omega$, the vector potential amplitude $\Psi$ of the photon must increase synchronously with the frequency ($\Psi \propto \omega$). This causes the energy density $u$ at the core of the photon to increase with the fourth power of the frequency: $u \propto \omega^4 g$.

\vspace{0.5em}
\noindent \textbf{2. Modification of the Scale Equation:} \\
Substituting $u \propto \omega^4$ into the logarithmic function $g$, we obtain the transcendental equation for the soliton equilibrium condition:

\begin{equation}
g = \kappa \ln(1 + C \omega^4 g) \approx 4\kappa \ln(\omega) 
\end{equation}

Solving this yields the transverse entity radius $R$ of the photon as:

\begin{equation}
R = \frac{a}{\sqrt{g}} = \frac{\lambda}{2\pi \sqrt{4\kappa \ln(\omega)}} 
\end{equation}

\vspace{0.5em}
\noindent \textbf{Analytical Conclusions:}

\begin{enumerate}
    \item \textbf{Fundamental Linear Relationship:} On a macroscopic scale, $R \propto \lambda$, which explains the geometric criterion in classical optics where wavelength determines resolution.
    \item \textbf{Logarithmic Scaling Violation:} Due to the presence of the $\sqrt{\ln \omega}$ term in the denominator, as the frequency increases, the surge in energy density leads to stronger nonlinear suppression by the grid. This means that high-energy photons (such as gamma rays) will be more compact and denser than predicted by a purely linear proportion.
    \item \textbf{Frequency-Adaptive Equilibrium:} By dynamically adjusting the smoothing scale $a=c/\omega$, the grid ensures that regardless of whether the frequency is high or low, the photon can achieve a balance between transverse diffraction and energy density self-focusing within its intrinsic cycle. This signifies that the photon is a non-collapsing, \textbf{three-dimensional topological soliton} whose scale adaptively adjusts to its frequency.
\end{enumerate}






\section{The Fluid Topological Origin of the Dirac Spinor and the Vortex Gap Soliton Solution}

In Section 11, we established the underlying grid's non-local saturated nonlinear field equation and verified the rationality of the photon as a spatial soliton solution. However, to thoroughly answer the ultimate question of ``what the electron entity exactly is,'' we must seek for it, within the four-dimensional spacetime manifold, an exact analytical topological state that simultaneously satisfies rest mass, spin $1/2$ ($4\pi$ symmetry), and magnetic moment properties. This section, combining the biorthogonal spacetime geometric picture of Clifford Algebra, will demonstrate that the \textbf{Dirac Equation} in relativistic quantum mechanics and its ``spinor'' structure are essentially the emergence of a \textbf{spatiotemporal vortex gap soliton} from the underlying nonlinear fluid equations.

\subsection{Physical Geometric Picture: Polarization-Phase Synchronization and Spatiotemporal Dual Spin}

According to the four-dimensional partition of spacetime by Clifford Algebra, the ``self-trapped standing wave'' of the free electron cannot be viewed as a simple three-dimensional spherically symmetric scalar oscillation, but must contain two mutually orthogonal rotational subspaces. We define 2 orthogonal subspaces:

\begin{enumerate}
    \item \textbf{Spatial orbital vortex in the $x-y$ plane:} The photon momentum flow undergoes circular closed motion within the transverse cross-section, generating intrinsic topological vortex angular momentum (this is the underlying source of the electron's macroscopic angular momentum and magnetic dipole moment).
    \item \textbf{Spatiotemporal wave oscillation in the $z-t$ plane:} The inherent wave nature of the photon continuously alternates and advances longitudinally over time, forming a longitudinal standing wave.
\end{enumerate}

To prevent the momentum flow from escaping along the $z$-axis, thereby forming a stable physical particle, the system must achieve a dynamical locking, which we term ``Polarization-Phase Synchronization'':\\
The electron is essentially an \textbf{ultra-strongly polarized photon}, whose transverse polarization axis rotational frequency (intrinsic polarization frequency) is strictly locked to its forward wave carrier frequency $\omega$.

Based on this physical picture, we construct an Ansatz for the complex vector transverse component $A_+ = A_x + iA_y$ of the electromagnetic four-potential $A^\mu$ in the cylindrical coordinate system $(\rho, \theta, z)$. Considering that the standing wave is formed by the superposition of forward and backward propagating waves, we set its morphology as:

\begin{equation}
A_+(t, \rho, \theta, z) = \left[ F(\rho, z) e^{ikz} + B(\rho, z) e^{-ikz} \right] e^{i(\theta - \omega t)} 
\end{equation}

Where:
\begin{itemize}
    \item $e^{-i\omega t}$ represents the circular polarization spin of the polarization axis over time;
    \item $e^{i\theta}$ represents the spatial topological vortex of the momentum flow in the $x-y$ plane;
    \item $e^{\pm ikz}$ represents the spatiotemporal traveling carrier wave along the $z$-axis (by the speed of light condition, wavenumber $k = \omega/c$);
    \item $F(\rho, z)$ and $B(\rho, z)$ are the slowly varying envelopes to be determined.
\end{itemize}

\subsection{The Ontology of Energy Density: Vector Cancellation and Scalar Pressure Standing Wave}

Before substituting the solution into the field equation, let us review the calculation rules for the grid energy density $u_{inv}$. As mentioned earlier, the grid energy density corresponds to the \textbf{scalar pressure superposition of the potential}, rather than the electromagnetic field strength tensor. When approaching the geometric center $\rho \to 0$, due to the symmetry of the topological vortex phase, the vector sum of the spatial transverse momentum vector flows cancels out to 0 (ensuring that the spatial fluid does not undergo multi-valued tearing, i.e., the central field strength $\vec{E}=0$). However, the scalar potential (representing the absolute momentum crowdedness/vacuum pressure of the grid) perfectly superimposes at the geometric center, reaching a local maximum. The electric field strength is the negative gradient of the electric potential; when the potential has an extremum at the center, the field strength tends to 0.

For the transverse vector potential itself, we calculate the square of its modulus $\mathbf{A}_\perp^2 = |A_+|^2$ to extract its spatial structure:

\begin{equation}
\mathbf{A}_\perp^2 \propto \Big| [F e^{ikz} + B e^{-ikz}] e^{i(\theta - \omega t)} \Big|^2 
\end{equation}
\begin{equation}
= \Big( [F e^{ikz} + B e^{-ikz}] e^{i(\theta - \omega t)} \Big) \cdot \Big( [F^* e^{-ikz} + B^* e^{ikz}] e^{-i(\theta - \omega t)} \Big) 
\end{equation}

This is a profound mathematical fact: although the potential field itself spins and revolves at extremely high frequencies (incorporating the dynamic high-frequency phases of $t$ and $\theta$), \textbf{these dynamic phases are strictly conjugated and canceled out in the product ($e^{i(\theta - \omega t)} \cdot e^{-i(\theta - \omega t)} \equiv 1$)}.\\
The expanded pure spatial static distribution is:

\begin{equation}
\mathbf{A}_\perp^2 = |F|^2 + |B|^2 + F B^* e^{2ikz} + F^* B e^{-2ikz} 
\end{equation}

\begin{itemize}
    \item \textbf{DC Background:} $|F|^2 + |B|^2$ constitutes the constant ``stationary pressure background'' of the vacuum grid.
    \item \textbf{Self-induced Bragg Grating (AC):} The cross-interference term $F B^* e^{2ikz} + c.c.$ does not contain time $t$. This indicates that the high-frequency dual-spinning photon flow, through its own interference, creates a stationary ``refractive index grating'' along the $z$-axis in the vacuum.
\end{itemize}

\subsection{Dynamic-Static Collapse of the D'Alembert Operator: Transition from Second-Order to First-Order}

Now, we substitute the Ansatz into the four-dimensional spacetime D'Alembert operator $\Box = \frac{1}{c^2}\partial_t^2 - \nabla_\perp^2 - \partial_z^2$.

\begin{enumerate}
    \item \textbf{Time Term:} $\frac{1}{c^2}\partial_t^2 A_+ = -\frac{\omega^2}{c^2} A_+ = -k^2 A_+$
    \item \textbf{Longitudinal Spatial Term:} Applying the product rule, $\partial_z^2 (F e^{ikz}) = (\partial_z^2 F + 2ik \frac{\partial F}{\partial z} - k^2 F) e^{ikz}$.
\end{enumerate}

Since the spatial variations of envelopes $F$ and $B$ are much slower than the carrier wavelength, we introduce the \textbf{Slowly Varying Envelope Approximation (SVEA)}, ignoring the tiny second-order derivative $\partial_z^2 F$.\\
Rearranging yields: the $-k^2$ representing ``forward momentum'' and the time eigenvalue $-k^2$ \textbf{exactly undergo a strict positive-negative cancellation}. The result after the operator action collapses to:

\begin{equation}
\Box A_+ \approx \left[ (-\nabla_\perp^2 F - 2ik\frac{\partial F}{\partial z}) e^{ikz} + (-\nabla_\perp^2 B + 2ik\frac{\partial B}{\partial z}) e^{-ikz} \right] e^{i(\theta - \omega t)} 
\end{equation}

Due to the intrinsic synchronization of the polarization frequency and the wave frequency, the D'Alembert operator, which originally contains second-order derivatives of time and space (a characteristic of boson propagation), is strictly reduced to an operator containing only \textbf{first-order spatial derivatives ($\partial_z$)}. This is the underlying pure mathematical-mechanical mechanism for the transition from a massless scalar field to a massive fermionic spinor field.

\subsection{Emergence of the Dirac Equation and the Pure Fluid Essence of Mass}

Substituting the dimensionally-reduced operator into the first-order approximation form $(1 + a^2\Box) A_+ = [1 + N(u_{inv})] A_+$ of the non-local saturated field equation from Section 11.4. The nonlinear potential energy term $N(u_{inv})$ separates into a DC term $N_{DC}$ and a grating reflection term $N_{Bragg}$.\\
By extracting the coefficients for $e^{ikz}$ and $e^{-ikz}$, we obtain a set of symmetric Coupled-Mode Equations (CMEs):

\begin{equation}
\begin{cases}
- i (2a^2 k) \frac{\partial F}{\partial z} - a^2 \nabla_\perp^2 F = N_{DC} F + N_{Bragg} B \\
+ i (2a^2 k) \frac{\partial B}{\partial z} - a^2 \nabla_\perp^2 B = N_{DC} B + N_{Bragg}^* F
\end{cases}
\end{equation}

\textbf{Macroscopic Isomorphism with the Dirac Equation:}\\
This set of macroscopic continuity equations, derived entirely from classical nonlinear medium fluid mechanics, is mathematically \textbf{isomorphic to the $(1+1)D$ Dirac Equation containing transverse momentum in relativistic quantum mechanics.}

\begin{itemize}
    \item The forward envelope $F$ and backward envelope $B$ perfectly correspond to the left-handed and right-handed components of the Dirac Bi-spinor.
    \item The first-order spatial derivative term $\pm i \partial_z$ strictly corresponds to the Dirac momentum operator $\boldsymbol{\alpha} \cdot \hat{p}$.
\end{itemize}

In the Dirac equation, the coefficient coupling the left-handed and right-handed spinors is exactly the rest mass term $mc^2$. In our fluid model, it strictly corresponds to the cross-coupling term coefficient $N_{Bragg}$.\\
\textbf{Physical Conclusion:} In our equation, the rest mass of the electron is \textbf{the reflection coupling coefficient of the ``self-induced Bragg grating'' carved into the vacuum grid by the photon under extreme energy density due to polarization-phase synchronization.} The reason the momentum flow can be stationary in space and exhibit macroscopic inertia is that, within this grating of its own making, it continuously undergoes an infinite cycle where the forward wave is reflected into the backward wave, and the backward wave is reflected into the forward wave. This dynamical process corresponds to the electron ``jitter'' (Zitterbewegung) described by Schr\"{o}dinger.

\subsection{Time-Dependent Solutions and Time-Dependent Coupled-Mode Equations: The Mechanical Origin of Zitterbewegung}

To comprehensively describe the characteristics of the electron as a dynamical entity and to answer its intrinsic connection with the \textbf{Zitterbewegung} of electrons in relativistic quantum mechanics, we extend the model to the time domain.

\subsubsection{Dynamic Envelope Ansatz and Operator Evolution}

We extend the static envelopes $F(\rho, z), B(\rho, z)$ from Section 12.1 into dynamic envelopes $F(z, t)$ and $B(z, t)$ that vary slowly over time. Considering the intrinsic frequency locking exhibited by the electron on a macroscopic scale, the Ansatz is modified to:

\begin{equation}
A_+(z, t) = \left[ F(z, t) e^{ikz} + B(z, t) e^{-ikz} \right] e^{-i\omega t} 
\end{equation}

Substitute this dynamic solution into the four-dimensional D'Alembert operator $\Box = \frac{1}{c^2}\partial_t^2 - \partial_z^2$ (temporarily ignoring the transverse $\nabla_\perp^2$ component to focus on longitudinal dynamics). Applying the Slowly Varying Envelope Approximation (SVEA), i.e., ignoring the second-order spatiotemporal derivatives of the envelope terms $\partial_t^2 F \approx 0, \partial_z^2 F \approx 0$.

The operator acting on the forward wave component $F e^{ikz} e^{-i\omega t}$ yields:

\begin{equation}
\Box \left( F e^{i(kz-\omega t)} \right) \approx \left[ -\frac{\omega^2}{c^2} F - \frac{2i\omega}{c^2} \frac{\partial F}{\partial t} - (ik)^2 F - 2ik \frac{\partial F}{\partial z} \right] e^{i(kz-\omega t)} 
\end{equation}

Using the synchronous locking condition of standing wave self-trapping $k^2 = \omega^2/c^2$, the second-order eigen terms (the basis of the mass term) once again achieve strict cancellation. After eliminating the canceled terms, the operator collapses into a linear combination containing first-order partial derivatives of time and space:

\begin{equation}
\Box \left( F \right) \approx -2ik \left( \frac{1}{c} \frac{\partial F}{\partial t} + \frac{\partial F}{\partial z} \right) e^{i(kz-\omega t)} 
\end{equation}

\subsubsection{Time-Dependent CMEs and Dirac Isomorphism}

Substitute the dynamic collapsed operator into the nonlinear field equation $a^2\Box A_+ = N_{Bragg} A_+$, and according to the mode-matching principle, extract the coefficient terms of $e^{ikz}$ and $e^{-ikz}$. Introducing the coupling constant $\kappa = -N_{Bragg}/(2a^2 k)$, we obtain the \textbf{Time-dependent Coupled-Mode Equations (Time-dependent CMEs)} describing the fluid envelope evolution:

\begin{equation}
\begin{cases}
i \left( \frac{\partial F}{\partial t} + c \frac{\partial F}{\partial z} \right) = (c\kappa) B \\
i \left( \frac{\partial B}{\partial t} - c \frac{\partial B}{\partial z} \right) = (c\kappa) F
\end{cases}
\end{equation}

Comparing this with the general form of the spinor component expansion of the standard $(1+1)D$ Dirac equation, the fluid envelopes $F, B$ and the Dirac spinors $\psi_1, \psi_2$ are completely equivalent in structure. The physical anchoring relationship between the two theoretical systems is given by the coupling term: \textbf{$c\kappa \equiv mc^2/\hbar$}.

In the Dirac theory, a free electron, even in a macroscopically resting state, possesses a high-frequency Zitterbewegung at an angular frequency $\omega_{ZB} = 2mc^2/\hbar$. In this fluid model, this ``jitter'' manifests as the spontaneous high-frequency beating of energy between the forward and backward flows when the particle is at rest ($\partial_z F, \partial_z B = 0$). For the envelope field, the \textbf{Beat Frequency} of its amplitude cyclic conversion is the difference between the two eigenfrequencies:

\begin{equation}
\omega_{beat} = (+c\kappa) - (-c\kappa) = 2c\kappa = \frac{2mc^2}{\hbar} = \omega_{ZB}
\end{equation}

\subsubsection{Summary}

The beat frequency of the fluid model is \textbf{absolutely identical} to the jitter frequency calculated by quantum mechanics. It is precisely described as \textbf{an entity physical process where the photon flow is continuously reflected at the speed of light within the self-induced Bragg grating, causing the forward momentum and backward momentum to periodically and rapidly exchange (beat frequency phenomenon).} The characteristic exchange length (i.e., the grating penetration depth) given by this model is $L = 1/\kappa = \hbar/mc$, which perfectly matches the reduced Compton wavelength and the Zitterbewegung amplitude scale of the electron.

\subsection{The Quantized Rupture of Speed: Why Doesn't a $0.5c$ Photon Exist?}

Above, we proved that photons can generate self-induced Bragg gratings through polarization rotation. An intuitive question follows: since spin can generate reflection pressure, does there theoretically exist a ``intermediate state''---i.e., a ``slow photon'' whose effective propagation speed caused by rotation is between $0$ and $c$ (e.g., $0.5c$)? Observational facts indicate that such an intermediate state does not exist in a vacuum. Under this framework, this is not accidental, but rather a stability selection caused by the \textbf{Phase-space Attractor} effect of the underlying nonlinear field equation.

\subsubsection{Two Topological Attractors in the Vacuum Phase Space}

According to ``Computational Realism,'' the vacuum is not a simple background, but a nonlinear dynamical system with a saturation threshold. In this system, the evolution of momentum flow has two extremely stable ``energy-minimized'' topological solutions, namely phase-space attractors:

\begin{enumerate}
    \item \textbf{Linear Propagator (Photon State):} Momentum translates at the absolute speed $c$. In this state, nonlinear effects are extremely weak, and the momentum flow maintains a delicate balance between transverse divergence and self-focusing.
    \item \textbf{Nonlinear Standing-Wave Soliton (Electron State):} Through the complete synchronization (Phase-locking) of polarization frequency and wave frequency, the momentum is locked dead in the self-induced grating, and the macroscopic translation speed is $0$ (when at rest).
\end{enumerate}

Between these two attractors, there exists a massive \textbf{Instability Gap}. We argue that any attempt to occupy an intermediate state between the electron state and the photon state will lead to a rapid disintegration of the topological structure because it cannot satisfy the ``perfect standing wave condition.''

\subsubsection{Low-Energy Limit: Trivial Superposition in the Maxwell Linear Domain}

When the photon's polarization rotation frequency is low (energy density far below the vacuum saturation threshold $u_{th}$), the vacuum exhibits excellent linear characteristics.

At this point, a circularly polarized photon is mathematically strictly equivalent to the linear superposition of two spatially orthogonal, linearly polarized photons with a phase difference of $\pi/2$. Because the nonlinear term $g(u) \approx 0$, the vacuum does not generate enough reflection pressure to alter the speed of light. The momentum flow does not feel the existence of a ``self-induced grating,'' so the wave packet always propagates at the vacuum intrinsic speed $c$. This is the domain where classical electrodynamics (Maxwell's equations) applies.

\subsubsection{High-Energy Limit: Polarization Instability and Orthogonal Splitting}

When the photon's energy density approaches the critical threshold, the vacuum's nonlinear coupling term $g(u)$ is violently activated. At this time, the circularly polarized state is no longer a stable analytical solution, but will undergo \textbf{Spontaneous Polarization Splitting}:

\begin{enumerate}
    \item \textbf{Nonlinear Birefringence:} Extremely high-energy circularly polarized photon flows will induce non-uniform refractive index gradients on the grid due to their own energy density distribution.
    \item \textbf{Mode Competition:} In nonlinear dynamics, circularly polarized states are often unstable extrema. Once disturbed, it will quickly split into two linearly polarized modes with lower energy and mutually perpendicular polarizations.
    \item \textbf{Energy Release:} This splitting is accompanied by the collapse of the topological structure, and the photon cannot maintain a uniform deceleration. It either maintains propagation at $c$ (by splitting or adjusting through frequency shift) or directly collapses into the potential well of the electron state through ``strong coupling.''
\end{enumerate}

\subsubsection{Convergence with Predictions of Modern Quantum Electrodynamics (QED)}

The conclusions of this theory are consistent with the cutting-edge understandings of high-energy photon instabilities in QED:

\begin{itemize}
    \item \textbf{Vacuum Birefringence:} The vacuum is no longer isotropic under strong magnetic fields or high energy densities, which is precisely the macroscopic manifestation of ``momentum pressure modulating the effective speed of light'' in this theory.
    \item \textbf{Vacuum Photon Splitting:} Under extremely strong fields, a high-energy photon will spontaneously split into two lower-energy photons. In this framework, this is exactly the energy release mechanism caused by the ``instability of intermediate speed states.''
\end{itemize}

\subsubsection{Summary}

Photons at $0.5c$ do not exist in nature because speed itself is an emergence of topological closure conditions. Photons and electrons are two topological stable solutions of the vacuum equations under ``zero feedback'' and ``full feedback,'' respectively. Intermediate states are dynamically unstable due to their inability to form a closed intrinsic phase chain. This explains the constancy of the speed of light and reveals \textbf{the mechanical origin of the quantized jump between matter (electron) and energy (photon).}

\subsection{Conclusion: Topological Biorthogonal Locking and $4\pi$ Symmetry}

Thus far, the complete topological structure of the electron under this framework has been established: it is a \textbf{Spatiotemporal Vortex Gap Soliton} emerging from the underlying information grid.\\
In the transverse $x-y$ direction, the topological vortex centrifugal force balances the nonlinear pressure self-focusing; in the longitudinal $z-t$ direction, the pressure Bragg grating carved by polarization synchronization locks dead the longitudinal escape path.\\
This biorthogonal phase locking requires the system not only to rotate $360^\circ$ in three-dimensional space (to restore the spatial vortex phase) but also to advance by a corresponding half-cycle in the time coordinate (because the spatial rotation induces a half-period evolution of the internal standing wave on the time manifold), so that both the forward and backward coupled wave packets jointly return perfectly to the exact initial state. This provides a classical fluid dynamics and topology proof for the $4\pi$ ($720^\circ$) geometric symmetry of fermion half-integer spin, thereby unifying classical field theory, nonlinear soliton dynamics, and relativistic quantum spinor structure within the same computational framework.


\end{document}
